%!TEX root = main.tex
%=========================================================

\begin{abstract}
Mobile Ad-Hoc Networks (MANETs) allow distributed applications where no fixed network infrastructure is available.
MANETs use wireless communication subject to faults and uncertainty, and must support efficient \emph{broadcast}.
Controlled flooding is suitable for highly-dynamic networks, while overlay-based broadcast is suitable for dense and more static ones.
Density and mobility vary significantly over a MANET deployment area.
We present the design and implementation of \emph{emergent overlays} for efficient and reliable broadcast in heterogeneous MANETs.
This adaptation technique allows nodes to automatically switch from controlled flooding to the use of an overlay.
Interoperability protocols support the integration of both protocols in a single heterogeneous system.
Coordinated adaptation policies allow regions of nodes to autonomously and collectively emerge and dissolve overlays.
Our simulation of the full network stack of 600 mobile nodes shows that emergent overlays reduce energy consumption, and improve reliability and coverage compared to single protocols and to two previously-proposed adaptation techniques.
\end{abstract}
